\begin{frame}{이 절의 핵심 세 가지}
  \begin{block}{SVM = 같은 문제, 다른 원칙}
    로지스틱회귀·GDA와 \textbf{같은 문제}(선형 결정 경계)를
    ``\textbf{마진을 최대화한다}''는 다른 원칙으로 풀고 —
    \textbf{원칙이 다르면 얻는 성질(서포트 벡터, 커널 트릭)도
    달라진다.}
  \end{block}
  \vskip0.4em
  \begin{enumerate}
    \item \kb{고차원 = ``휘어진 경계''를 ``평평한 경계''로}:
      원형 데이터는 $z=x_1^2+x_2^2$ 축을 추가하면 두 수평 고도로
      벌어져 평면 하나로 잘림. XOR은 $x_1x_2$ 축이 필요.
    \item \kb{커널 트릭 = ``내적만 계산하자'':} 쌍대 문제에서
      데이터는 $\phi(x^{(i)})\cdot\phi(x^{(j)})$ 로만 등장하므로,
      이 내적값을 $K(x,x')$ 으로 교체해도 해가 변하지 않음.
      2$\times$2 커널 행렬 예제로 primal과 dual이 같은 해
      ($w=(1/6,1/6)$)을 줌을 확인.
    \item \kb{$\gamma$(또는 다항 커널의 도수 $d$)는 ``경계의
      구불 복잡도'' 손잡이}: $\gamma$ 가 작으면 직선(과소적합),
      크면 점마다 봉우리(과적합). \textbf{반드시 $(C,\gamma)$ 를
      함께 검증셋으로 튜닝}.
  \end{enumerate}
  \vskip0.6em
  \begin{block}{6장(정규화와 모델 선택)으로 넘어가기 전, 연결고리}
    이번 절의 $C$, $\gamma$ 튜닝은 6장의 \textbf{교차 검증}으로
    체계화 — 검증셋 정확도로 $(C,\gamma)$ 를 함께 그리드 서치가
    표준 관행.
    10장(신경망)에서 RBF의 ``봉우리''와 CNN의 \textbf{커널 필터}는
    수학적으로 다른 개념이지만, ``원본 특징을 고차 항으로 변환해
    선형 분류''한다는 큰 그림은 이어진다.
  \end{block}
\end{frame}
